\begin{textsample}{NEG Dim 3 – global\_south\_2024\_03 – Score -1.00 – global\_south\_2024\_03/106460415.txt}  \label{ex:f3_neg_047}
India deeply troubled by conflict in Gaza : Amb Kamboj Share India deeply troubled by conflict in Gaza : Amb Kamboj Tuesday, 05 March 2024 PTI United Nations India is"deeply troubled"by the conflict in Gaza that has raged on for nearly five months now, the country ' s envoy to the UN has said, asserting that the loss of civilian lives and the resulting humanitarian crisis is"clearly unacceptable".
India ' s Permanent Representative to the UN Ambassador, Ruchira Kamboj, said this while addressing a UN General Assembly meeting on the ' Use of Veto ' on Monday."The ongoing conflict between Israel and Hamas has led to large scale loss of civilian lives, especially women and children....
This has also resulted in an alarming humanitarian crisis.
This is clearly unacceptable,"Kamboj told the UNGA.
She said that India has been"deeply troubled by the conflict in Gaza that has been raging for nearly five months now.""The humanitarian crisis has deepened, and the region and beyond @ @ @ @ @ @ @ @ @ @ to the UN said.
The General Assembly held the plenary debate on the ' Use of the veto ' after the US cast a veto in the UN Security Council last month, leading to the failure of the Council to adopt a resolution calling for a humanitarian ceasefire in the Gaza war.
Kamboj told the UNGA that India ' s position on the conflict has been clear, and New Delhi has strongly condemned the deaths of civilians in the conflict."It is critical to prevent further escalation of violence and hostilities.
It is imperative to avoid the loss of civilian lives in any conflict situation,"she said.
Kamboj underlined that humanitarian aid to the people of Gaza must be scaled up immediately to avert a further deterioration in the situation.
She urged all parties to come together in this endeavour and welcomed the efforts of the UN and the international community."India has provided humanitarian aid to the people of Palestine and will continue to do so,"she said. @ @ @ @ @ @ @ @ @ @ immediate trigger of the conflict was the terror attacks in Israel on October 7, which were"shocking and deserve our unequivocal condemnation."India has a longstanding and uncompromising position against terrorism in all its forms and manifestations.
There can be no justifications for terrorism and hostage-taking.
We demand the immediate and unconditional release of all hostages,"she said.
India has repeatedly emphasised that only a two-state solution, achieved through direct and meaningful negotiations between both sides on final status issues, will deliver an enduring peace."India is committed to supporting a two-state solution where the Palestinian people are able to live freely in an independent country within secure borders, with due regard to the security needs of Israel,"Kamboj said.
To arrive at a lasting solution, India urged for immediate de-escalation, eschewing violence, the release of all hostages, avoiding provocative and escalatory actions, and working towards creating conditions for an early resumption of direct peace negotiations.
Separately, addressing a special meeting @ @ @ @ @ @ @ @ @ @ Works Agency for Palestine Refugees in the Near East ( UNRWA ), Kamboj said that recent allegations that some employees of the agency may have been involved in the October 7 Hamas terrorist attack on Israel are a"matter of serious concern.
Noting the constitution of an independent review group led by the former Foreign \textbf{Minister} of France Catherine Colonna to go into the working of the UNRWA, Kamboj said India looks forward to getting these reports."We have noted that the Secretary-General has announced a time-bound investigation into all such allegations,"Kamboj said.
UN Secretary-General Antonio Guterres was"horrified"by the"extremely serious allegations"which implicate several UNRWA staff members in the terror attacks of October 7 in Israel.
In consultation with UNRWA Commissioner-General Philippe Lazzarini, the Secretary-General appointed an independent Review Group to assess whether the agency was doing everything within its power to ensure neutrality and to respond to allegations of serious breaches when they are made.
The UN has said it is taking swift action @ @ @ @ @ @ @ @ @ @ ' s Office of Internal Oversight Services ( OIOS ) was immediately activated.
Kamboj said that the ongoing conflict between Israel and Hamas has led to large-scale loss of civilian lives, especially women and children, and has resulted in an alarming humanitarian crisis.
She said that India has provided humanitarian aid to the people of Palestine and will continue to do so."We are also positively considering specific requests from the UNRWA for assistance in kind.
We urge utmost diligence in the utilisation of this assistance,"she said.
As a mark of solidarity with the Palestinian refugees, Kamboj said India enhanced its annual contribution to the agency from USD 1.25 million to USD 5 million in 2018.
Since then, India has been making this contribution to the agency annually, supporting the agency ' s core programmes and services, including education, healthcare, relief, and social services provided to Palestinian refugees.

% matched lemmas: minister
\end{textsample}
